\documentclass[./../main.tex]{subfiles}

\begin{document}
\section{chapter2 归纳法}

\subsection{归纳算法举例}
\subsubsection{多项式求值}
假设有$n+2$个实数$a_0,a_1,...,a_n$和$x$的序列，现在要对以下多项式求值：
\[
    P_n(x)=\sum_{i=0}^{n}a_ix^i
\]
我们发现：
\[
    p_n(x)=((...(((a_nx+a_{n-1})x+a_{n-2})x+a_{n-3})...))x+a_0
\]
所以，我们得到具体表达式为：
\[
    p_{n-1}(x)=a_nx^{n-1}+a_{n-1}x^{n-2}+...+a_2x+a_1
\]
所以，有递推关系：
\[
    p_j(x)=xp_{j-1}(x)+a_{n-j}
\]

\begin{algorithm}
    \SetAlgoLined
    \KwIn{$n+2$个实数$a_0,a_1,...,a_n$和$x$的序列}
    \KwOut{表达式$P_n(x)$的值}
    \BlankLine
    $p\leftarrow a_n$\\
    \For{$j\leftarrow1$ \KwTo $n$}{
        $p\leftarrow xp+a_{n-j}$\\
    }
    \Return $p$
    \caption{HORNER}
\end{algorithm}

\subsubsection{生成排列}
\begin{algorithm}[h]
    \KwIn{正整数$n$}
    \KwOut{所有长度为$n$的排列}
    \BlankLine
    \For{$j\leftarrow1$ \KwTo $n$}{
        $P[j]\leftarrow j$
    }
    perm1(1)\\
    perm1(m)\\
    \eIf{$m=n$}{
        \Return P\tcp{output P}
    }{
        \For{$j\leftarrow m$ \KwTo $n$}{
            swap($P[j], P[m]$)\\
            perm1(m+1)\\
            swap($P[j], P[m]$)\tcp{$P[m..n]=m,m+1,...,n$}
        }
    }
    \caption{PERMUTATIONS1}
\end{algorithm}

显然，对于重复次数$f(n)$有以下递推关系：
\[
    f(n)=nf(n-1)+n, n\ge2
\]
令$f(n)=n!h(n)$, 那么
\[
    n!h(n)=n(n-1)!h(n-1)+n
\]
即
\[
    h(n)=h(n-1)+\frac{1}{(n-1)!}
\]
解得
\[
    h(n)=h(1)+\sum_{i=1}^{n-1}\frac{1}{i!}<\sum_{i=1}^{\infty}=e-1
\]
所以
\[
    f(n)=n!h(n)<(e-1)n!<2\cdot n!
\]
而输出语句的复杂度是$\Theta(n)$，所以最终$T(n)=\Theta(n\cdot n!)$

这里用到的思想其实是一种累加法的变体。也就是考虑递推式
\[
    f(n)=g(n)f(n-1)+h(n)
\]

只需进行合适的换元，令$p(n)$满足：
\[
    f(n)=(\prod_{i=1}^{n}g(i))\cdot p(n)
\]
这样得到：
\[
    (\prod_{i=1}^{n}g(i))\cdot p(n)=(\prod_{i=1}^{n}g(i))\cdot p(n-1)+ h(n)
\]
\[
    p(n)=p(n-1)+\frac{1}{\prod_{i=1}^{n}g(i)}h(n)
\]
\[
    p(n)=p(1)+\sum_{i=1}^{n}\frac{1}{\prod_{j=1}^{i}g(j)}h(i)
\]
由此，即可求出$f(n)$，不过要注意无穷级数的作用。

\subsubsection{基数排序}
令 $L = \{a_{1}, a_{2}, \ldots, a_{n}\}$ 是一张有 $n$ 个数的集合，其中每个数恰好有 $k$ 位数字，就是说每个数具有 $d_{k}d_{k-1}\ldots d_{1}$ 形式，这里 $d_{i}$ 是 $0$ 到 $9$ 中的一个数字。

我们研究一种排序算法对 $L$ 进行非降序排列，使算法在几乎所有的实际应用中都以线性时间运行。

在这个问题中，我们对整数的大小 $k$ 施行归纳，而不是对数 $n$ 这个对象来用归纳法。

\begin{algorithm}[t]
    \KwIn{一个有$n$个数的表$L=\{a_1,...,a_n\}$和$k$位数字}
    \KwOut{按非非降序排列的表$L$}
    \BlankLine
    \For{$j=1$ \KwTo $k$}{
        准备10个空表$L_0,L_1,...,L_9$
        \While{$L\neq\emptyset$}{
            $a\gets L.top()$\\
            $L.pop()$\\
            $i\leftarrow a$的第$j$位数字，将$a$加入表$L_i$中。\tcp{取个位为第1位数字}
        }
        $L\leftarrow L_0$\\
        \For{$i=1$ \KwTo $9$}{
            $L\leftarrow L\cup L_i$
        }
    }
    \Return{$L$}
    \caption{基数排序}
\end{algorithm}
时间复杂度：$\Theta(kn)$, 空间复杂度：$\Theta(10n)$

\subsubsection{寻找多数元素}
令 $A[1..n]$ 是一个整数序列，$A$ 中的整数 $a$ 如果在 $A$ 中出现的次数多于 $\lfloor n/2 \rfloor$，那么 $a$ 称为多数元素。


\begin{itemize}
    \item 初始，令候选元素$c=A[1]$，并将计数器置$1$
    \item 从$A[2]$开始，逐个扫描元素，如果被扫描的元素和$c$相等，则计数器加$1$；如果元素不等于$c$，则计数器减$1$
    \item 如果所有的元素都已经扫描完毕并且计数器大于$0$，那么返回$c$作为多数元素的候选者
    \item 如果在$c$和$A[j]$（$1<j<n$）比较时计数器为$0$，那么对于$A[j+1..n]$上的元素递归调用candidate过程。
\end{itemize}

\begin{algorithm}[h]
    \KwIn{一个数组 $A$ 和整数 $m, n$}
    \KwOut{数组 $A$ 中的一个候选者}
    \BlankLine
    $j \gets m$; $c \gets A[m]$; $count \gets 1$\;
    \While{$j < n \text{ and }count > 0$}{
        $j \gets j + 1$\;
        \eIf{$A[j] = c$}{
            $count \gets count + 1$\;
        }{
            $count \gets count - 1$\;
        }
    }
    \eIf{$j = n$}{
        \Return{$c$} \tcp{获取候选者}
    }{
        \Return{candidate($j+1$)}
    }
    \caption{candidate(m)}
\end{algorithm}

\begin{algorithm}[h]
    \KwIn{$n$ 个元素的数组 $A[1..n]$}
    \KwOut{若存在则输出多数元素，否则输出 none}
    \BlankLine
    $c \gets \text{candidate}(1)$\; \tcp{获取候选者}
    $count \gets 0$\;
    \For{$j \gets 1$ \KwTo $n$}{
        \If{$A[j] = c$}{
            $count \gets count + 1$\; \tcp{候选者进行计数}
        }
    }
    \If{$count > \lfloor n/2 \rfloor$}{
        \Return{$c$}\;
    }
    \Else{
        \Return{none}\;
    }
    \caption{MAJORITY}
\end{algorithm}

时间复杂度$\Theta(n)$，空间复杂度$\Theta(1)$

\subsubsection{棋子移动游戏}
洛谷题：\url{https://www.luogu.com.cn/problem/P1259}

对于最后处理的4对棋子是需要特殊对待的，不然会在2白2空2黑的时候卡死。

\subsection{本章习题}
\subsubsection{RADIXSORT算法时间复杂性}
当输入由下列区间中的$n$个正整数组成时，以$n$为大小说明算法RADIXSORT的时间复杂性。

\begin{enumerate}
    \item $[1 \ldots n]$
    \item $[1 \ldots n^{2}]$
    \item $[1 \ldots 2^{n}]$
\end{enumerate}

\begin{table}[ht]
\centering
\caption{基数排序（RADIXSORT）的时间复杂性分析}
\begin{tabular}{l >{$}l<{$} p{8cm}}
\toprule
\textbf{输入区间} & \textbf{时间复杂度} & \textbf{解释} \\
\midrule
$[1 \ldots n]$     & O(n)      & 最大位数$d = \lceil \log_{10} n \rceil \approx O(1)$，视为常数，故线性时间。 \\
$[1 \ldots n^2]$   & O(n^2)    & 位数$d = \lceil 2\log_{10} n \rceil$，但数字范围扩大，近似为$d = O(n)$。 \\
$[1 \ldots 2^n]$   & O(n \cdot 2^n) & 位数$d = \lceil n\log_{10} 2 \rceil = \Theta(n)$，且输入规模为$2^n$量级。 \\
\bottomrule
\end{tabular}
\end{table}

\subsubsection{寻找名人问题}
在一个房间里有$n$个人，其中一个是名人。所谓名人，就是大家都认识他，但是他却不认识任何其他人。名人之外的其他人可能认识房间里另外的一部分人。你可以问房间里任何人问题，但是问题只能是：“你认识某某吗？”，对方回答“Yes”或者“No”。请设计一个时间复杂度为$\Theta(n)$的归纳算法，通过问问题，把房间中的名人找出来。

\begin{enumerate}
    \item \textbf{初始化候选人}：假设第一个人为候选名人 \texttt{candidate = 0}。
    
    \item \textbf{遍历排除非名人}：
    
    \item 从第2个人到第$n$个人，依次与候选人比较：
    \begin{itemize}
        \item 若候选人认识当前人，则候选人不是名人，更新候选人为当前人。
        \item 若候选人不认识当前人，则当前人不是名人，保留候选人。
    \end{itemize}
    \item \textbf{验证候选人合法性}：
    \begin{itemize}
        \item 检查是否所有人认识候选人。
        \item 检查候选人是否不认识任何人。
    \end{itemize}
\end{enumerate}

\end{document}